\begin{tikzpicture}[
  >=Stealth,
  % ---------- Estilos ----------
  celda/.style={
    draw=neuronaborde!70, line width=0.8pt, fill=entrada!35
  },
  activa/.style={
    draw=neuronaborde!70, line width=0.8pt, fill=neurona!75
  },
  kernel/.style={
    draw=neuronaborde!70, line width=0.8pt, fill=perdida!70
  },
  marco/.style={
    draw=neuronaborde, line width=2pt
  },
  bloque/.style={
    draw=neuronaborde, line width=1.2pt
  },
  flecha/.style={
    -{Stealth[length=3mm,width=2.4mm]},
    line width=1.3pt, color=conexion!85!black
  },
  union/.style={
    line width=1.3pt, color=conexion!85!black
  },
  etiqueta/.style={
    text=conexion, font=\Large
  }
]
  % ---------- Colores (locales a este tikzpicture) ----------
  \definecolor{entrada}{RGB}{120,190,255}
  \definecolor{neurona}{RGB}{100,175,245}
  \definecolor{neuronaborde}{RGB}{25,80,150}
  \definecolor{perdida}{RGB}{150,205,250}
  \definecolor{conexion}{RGB}{90,100,120}

  % ================= Entrada I (2x2) =================
  % El elemento resaltado es el que se "esparce" en el resultado.
  \foreach \el/\ic/\ir/\hh in {a/0/0/1, b/1/0/0, c/0/1/0, d/1/1/0} {
    \ifnum\hh=1
      \node[activa, minimum size=1cm, inner sep=0pt, anchor=north west]
        at (\ic,-\ir) {\large $\el$};
    \else
      \node[celda, minimum size=1cm, inner sep=0pt, anchor=north west]
        at (\ic,-\ir) {\large $\el$};
    \fi
  }
  \draw[marco] (0,0) rectangle (1,-1);
  \node[etiqueta] at (1,-3.7) {$I$};

  % ================= N\'ucleo K (2x2) =================
  \foreach \ke/\kc/\kf in {j/0/0, k/1/0, l/0/1, m/1/1} {
    \node[kernel, minimum size=1cm, inner sep=0pt, anchor=north west]
      at (4.2+\kc,-\kf) {\large $\ke$};
  }
  \draw[marco] (4.2,0) rectangle (6.2,-2);
  \node[etiqueta] at (5.2,-3.7) {$K$};

  % ================= Operadores =================
  \node[etiqueta] at (3.1,-1) {{\Huge $\ast^{T}$}};
  \node[etiqueta] at (7.5,-1) {\Huge $=$};

  % ================= Resultado C (4x4 = 4 bloques de 2x2) =================
  % Cada elemento de I se multiplica por K completo y ocupa su bloque 2x2 (stride 2).
  \foreach \el/\qc/\qf/\hh in {a/0/0/1, b/1/0/0, c/0/1/0, d/1/1/0} {
    \foreach \ke/\kc/\kf in {j/0/0, k/1/0, l/0/1, m/1/1} {
      \ifnum\hh=1
        \node[activa, minimum width=1.2cm, minimum height=1cm, inner sep=0pt,
              anchor=north west]
          at ({9+1.2*(2*\qc+\kc)},{1-1.0*(2*\qf+\kf)}) {\large $\el\ke$};
      \else
        \node[celda, minimum width=1.2cm, minimum height=1cm, inner sep=0pt,
              anchor=north west]
          at ({9+1.2*(2*\qc+\kc)},{1-1.0*(2*\qf+\kf)}) {\large $\el\ke$};
      \fi
    }
  }
  % Bordes de los cuatro bloques; el resaltado es el que genera la flecha.
  \draw[bloque] (9,1)   rectangle (11.4,-1);
  \draw[bloque] (11.4,1) rectangle (13.8,-1);
  \draw[bloque] (9,-1)  rectangle (11.4,-3);
  \draw[bloque] (11.4,-1) rectangle (13.8,-3);
  \draw[marco]  (9,1)   rectangle (11.4,-1);
  \node[etiqueta] at (11.4,-3.7) {$C$};

  % ================= Flechas: de I y K al resultado =================
  \draw[flecha] (0.5,0) -- (0.5,0.5) -- (9,0.5);
  \draw[union]  (5.2,0) -- (5.2,0.5);
\end{tikzpicture}